    \expectedproblemsection{Memorization Problems}

    \begin{expectedproblem}
        The Clay Mathematics Institute's Yang--Mills Existence and Mass Gap
        problem asks for a nontrivial quantum Yang--Mills theory on
        $\R^d$ with a positive mass gap. Determine $d$, and state whether this
        case has been solved.
    \end{expectedproblem}

    \begin{solution}
        As recorded in Facts to Memorize, the Millennium Prize Problem is posed
        on four-dimensional Euclidean space:
        \[
            \boxed{d=4}.
        \]
        As of 2026, the required rigorous construction and proof of a positive
        mass gap in this dimension remain open.
    \end{solution}

    \begin{expectedproblem}
        John Milnor first exhibited exotic spheres in dimension $n$: smooth
        manifolds homeomorphic, but not diffeomorphic, to the standard sphere
        $S^n$. Later work showed that the group $\Theta_n$ of oriented
        diffeomorphism classes of homotopy $n$-spheres has exactly $k$
        elements in this dimension. Determine
        \[
            n+k.
        \]
    \end{expectedproblem}

    \begin{solution}
        Milnor's first exotic spheres were seven-dimensional, and
        Kervaire--Milnor theory gives
        \[
            \Theta_7\cong\Z/28\Z.
        \]
        Thus, with orientation taken into account, there are $28$ classes in
        total: one standard class and $27$ exotic classes. Therefore $n=7$
        and $k=28$, so
        \[
            \boxed{n+k=35}.
        \]
    \end{solution}

    \begin{expectedproblem}
        This public-key cryptosystem was developed in 1977 by Ronald Rivest,
        Adi Shamir, and Leonard Adleman. It supports public-key encryption and
        digital signatures, and its standard key construction uses a modulus
        formed from two large primes. The three researchers received the 2002
        ACM A. M. Turing Award for their contribution to practical public-key
        cryptography. Identify the cryptosystem.
    \end{expectedproblem}

    \begin{solution}
        The cryptosystem is
        \[
            \boxed{\text{RSA}}.
        \]
        Its name consists of the initials of Rivest, Shamir, and Adleman.
    \end{solution}

    \begin{expectedproblem}
        For a connected finite graph, Euler's criterion states that an Euler
        circuit exists exactly when every vertex has even degree, while an
        Euler trail with distinct endpoints exists exactly when precisely two
        vertices have a certain degree parity. What kind of vertices occur
        zero times in the first case and exactly twice in the second?
    \end{expectedproblem}

    \begin{solution}
        They are
        \[
            \boxed{\text{vertices of odd degree}}.
        \]
    \end{solution}

    \begin{expectedproblem}
        This Norwegian mathematician lived from $1802$ to $1829$ and died at
        the age of $26$. Despite poverty and illness, he made fundamental
        contributions to analysis and algebra. He is especially known for
        proving that the general quintic equation cannot be solved by radicals.
        Identify the mathematician.
    \end{expectedproblem}

    \begin{solution}
        The mathematician is
        \[
            \boxed{\text{Niels Henrik Abel}}.
        \]
    \end{solution}

    \begin{expectedproblem}
        This closed surface has no boundary and is nonorientable, so it has no
        globally consistent distinction between two sides. Unlike a
        M\"obius strip, it is closed. It cannot be embedded in
        three-dimensional Euclidean space, although it can be represented
        there by a self-intersecting immersion. Identify the surface.
    \end{expectedproblem}

    \begin{solution}
        The surface is the
        \[
            \boxed{\text{Klein bottle}}.
        \]
    \end{solution}

    \begin{expectedproblem}
        This German mathematician was born near K\"onigsberg in $1862$, held a
        professorship at the University of K\"onigsberg, and moved to the
        University of G\"ottingen in $1895$. His work shaped algebra, number
        theory, analysis, and the axiomatic foundations of geometry, and his
        \emph{Grundlagen der Geometrie} appeared in $1899$. Identify the
        mathematician.
    \end{expectedproblem}

    \begin{solution}
        The mathematician is
        \[
            \boxed{\text{David Hilbert}}.
        \]
    \end{solution}

    \begin{expectedproblem}
        In game theory, this is a strategy profile in which no player can
        improve their payoff by changing only their own strategy while all
        other players keep theirs fixed. It is named after the mathematician
        portrayed in the film \emph{A Beautiful Mind}. What is this concept?
    \end{expectedproblem}

    \begin{solution}
        The concept is a
        \[
            \boxed{\text{Nash equilibrium}}.
        \]
    \end{solution}

    \begin{expectedproblem}
        A complex number is called this if it is not a root of any nonzero
        polynomial with rational coefficients. The numbers $\pi$ and $e$ are
        standard examples. What is such a number called?
    \end{expectedproblem}

    \begin{solution}
        It is called a
        \[
            \boxed{\text{transcendental number}}.
        \]
    \end{solution}

    \begin{expectedproblem}
        A positive integer is called this if it equals the sum of its positive
        proper divisors. For example,
        \[
            6=1+2+3,
            \qquad
            28=1+2+4+7+14.
        \]
        The Euclid--Euler theorem states that every even integer of this type
        has the form
        \[
            2^{p-1}(2^p-1),
        \]
        where $2^p-1$ is prime. What is such an integer called?
    \end{expectedproblem}

    \begin{solution}
        It is called a
        \[
            \boxed{\text{perfect number}}.
        \]
    \end{solution}

    \begin{expectedproblem}
        An Euler path traverses every edge of a graph exactly once. What is the
        standard name for a path that visits every vertex of a graph exactly
        once?
    \end{expectedproblem}

    \begin{solution}
        It is called a
        \[
            \boxed{\text{Hamiltonian path}}.
        \]
    \end{solution}

    \begin{expectedproblem}
        The obverse of the Fields Medal depicts a mathematician from antiquity,
        while the reverse depicts a sphere inscribed in a cylinder. Identify the
        mathematician and state the volume ratio of the sphere to the
        circumscribed cylinder.
    \end{expectedproblem}

    \begin{solution}
        The mathematician is Archimedes. If the sphere has radius $r$, then the
        cylinder has radius $r$ and height $2r$, so
        \[
            \frac{V_{\mathrm{sphere}}}{V_{\mathrm{cylinder}}}
            =
            \frac{\frac43\pi r^3}{2\pi r^3}
            =
            \frac23.
        \]
        Thus the answer is
        \[
            \boxed{\text{Archimedes and }2:3}.
        \]
    \end{solution}

    \begin{expectedproblem}
        Match the four 2018 Fields Medalists with the following principal
        mathematical anchors:
        \begin{problemchoices}
            \item $p$-adic geometry;
            \item boundedness of Fano varieties and the BAB conjecture;
            \item optimal transport;
            \item number theory and homogeneous dynamics.
        \end{problemchoices}
    \end{expectedproblem}

    \begin{solution}
        The Fields Medal table and laureate anchors in Facts to Memorize give the
        matching
        \[
            \boxed{
                \begin{aligned}
                    \text{Peter Scholze}
                    &\longleftrightarrow p\text{-adic geometry},\\
                    \text{Caucher Birkar}
                    &\longleftrightarrow \text{Fano varieties and BAB},\\
                    \text{Alessio Figalli}
                    &\longleftrightarrow \text{optimal transport},\\
                    \text{Akshay Venkatesh}
                    &\longleftrightarrow
                    \text{number theory and homogeneous dynamics}.
                \end{aligned}
            }
        \]
    \end{solution}

    \begin{expectedproblem}
        Maryna Viazovska proved optimal sphere-packing results in dimensions
        $8$ and $24$. What is the classical three-dimensional sphere-packing
        conjecture called, and who proved it?
    \end{expectedproblem}

    \begin{solution}
        The three-dimensional statement is the
        \[
            \boxed{\text{Kepler conjecture}},
        \]
        proved by
        \[
            \boxed{\text{Thomas Hales}}.
        \]
        The later Flyspeck project formally verified the proof.
    \end{solution}

    \begin{expectedproblem}
        Which of the following is not one of the combinatorial log-concavity
        conjectures associated with June Huh's work?
        \begin{problemchoices}
            \item Rota's conjecture;
            \item Mason's conjecture;
            \item Dowling--Wilson conjecture;
            \item Sullivan conjecture.
        \end{problemchoices}
    \end{expectedproblem}

    \begin{solution}
        The first three belong to the circle of matroid and combinatorial
        log-concavity results associated with June Huh and his collaborators.
        The Sullivan conjecture belongs to algebraic topology. Therefore, the
        odd one out is
        \[
            \boxed{\text{Sullivan conjecture}}.
        \]
    \end{solution}

    \begin{expectedproblem}
        Identify the Abel Prize laureate or laureates for each year:
        \[
            2018,\quad 2019,\quad 2020,\quad 2021,\quad 2022.
        \]
    \end{expectedproblem}

    \begin{solution}
        The sequence is
        \[
            \boxed{
                \begin{aligned}
                    2018&:\ \text{Robert Langlands},\\
                    2019&:\ \text{Karen Uhlenbeck},\\
                    2020&:\ \text{Hillel Furstenberg and Gregory Margulis},\\
                    2021&:\ \text{L\'aszl\'o Lov\'asz and Avi Wigderson},\\
                    2022&:\ \text{Dennis Sullivan}.
                \end{aligned}
            }
        \]
    \end{solution}

    \begin{expectedproblem}
        Give the Fields Medalists conventionally associated with the following
        national firsts:
        \[
            \text{Brazil},
            \qquad
            \text{Japan},
            \qquad
            \text{Iran},
            \qquad
            \text{Vietnam}.
        \]
    \end{expectedproblem}

    \begin{solution}
        The standard associations are
        \[
            \boxed{
                \begin{aligned}
                    \text{Brazil}
                    &\longleftrightarrow \text{Artur Avila},\\
                    \text{Japan}
                    &\longleftrightarrow \text{Kunihiko Kodaira},\\
                    \text{Iran}
                    &\longleftrightarrow \text{Maryam Mirzakhani},\\
                    \text{Vietnam}
                    &\longleftrightarrow \text{\NgoBaoChau}.
                \end{aligned}
            }
        \]
    \end{solution}




    \begin{expectedproblem}
        Which of the following mathematicians was born earliest?
        \begin{problemchoices}
            \item Guillaume de l'H\^opital;
            \item Isaac Newton;
            \item Gottfried Wilhelm Leibniz;
            \item Colin Maclaurin.
        \end{problemchoices}
    \end{expectedproblem}

    \begin{solution}
        Newton was born in $1643$ under the modern Gregorian calendar
        convention, earlier than Leibniz in $1646$, l'H\^opital in $1661$, and
        Maclaurin in $1698$. Therefore the answer is
        \[
            \boxed{\text{Isaac Newton}}.
        \]
    \end{solution}

    \begin{expectedproblem}
        What is Hilbert's first problem in the list of twenty-three problems
        presented in $1900$?
    \end{expectedproblem}

    \begin{solution}
        Hilbert's first problem asks for the resolution of the
        \[
            \boxed{\text{continuum hypothesis}}.
        \]
        In modern set theory, Gödel and Cohen showed that the continuum
        hypothesis is independent of ZFC, assuming ZFC is consistent.
    \end{solution}

    \begin{expectedproblem}
        Match each number system with the mathematician most directly
        associated with its introduction or early systematic development. The
        mathematicians are Hamilton, Hensel, Hewitt, and Conway.
        \begin{problemchoices}
            \item quaternions;
            \item $p$-adic numbers;
            \item hyperreal numbers;
            \item surreal numbers.
        \end{problemchoices}
    \end{expectedproblem}

    \begin{solution}
        The standard associations are
        \[
            \boxed{
                \begin{aligned}
                    \text{quaternions}
                    &\longleftrightarrow \text{Hamilton},\\
                    p\text{-adic numbers}
                    &\longleftrightarrow \text{Hensel},\\
                    \text{hyperreal numbers}
                    &\longleftrightarrow \text{Hewitt},\\
                    \text{surreal numbers}
                    &\longleftrightarrow \text{Conway}.
                \end{aligned}
            }
        \]
        Robinson later made hyperreal methods central to nonstandard analysis,
        and Knuth introduced the name \emph{surreal numbers}.
    \end{solution}

    \begin{expectedproblem}
        Match each clue with the corresponding paradox.
        \begin{problemsubparts}
            \item Among $23$ people, a shared birthday has probability greater
                  than $1/2$.
            \item Alternating two losing games can create a winning game.
            \item A ball can be reassembled into two congruent copies using
                  finitely many nonmeasurable pieces.
            \item Dominance reasoning and expected utility reasoning recommend
                  different choices when a highly accurate predictor has already
                  filled two boxes.
        \end{problemsubparts}
    \end{expectedproblem}

    \begin{solution}
        The matching is
        \[
            \boxed{
                \begin{aligned}
                    \text{(a)}\;&\text{birthday paradox},\\
                    \text{(b)}\;&\text{Parrondo's paradox},\\
                    \text{(c)}\;&\text{Banach--Tarski paradox},\\
                    \text{(d)}\;&\text{Newcomb's paradox}.
                \end{aligned}
            }
        \]
    \end{solution}

    \begin{expectedproblem}
        Euclid's fifth postulate is often replaced in modern elementary
        geometry by the statement that through a point outside a given line
        there is exactly one parallel line. What is this equivalent statement
        called?
    \end{expectedproblem}

    \begin{solution}
        It is
        \[
            \boxed{\text{Playfair's axiom}}.
        \]
    \end{solution}

    \begin{expectedproblem}
        Match the three constant-curvature geometries with the number of
        parallels through a point outside a given line and with the angle sum of
        a triangle.
        \begin{problemsubparts}
            \item Euclidean geometry;
            \item hyperbolic geometry;
            \item elliptic or spherical geometry.
        \end{problemsubparts}
    \end{expectedproblem}

    \begin{solution}
        The matching is
        \[
            \boxed{
                \begin{aligned}
                    \text{(a) Euclidean:}&\quad
                    \text{exactly one parallel, angle sum }=\pi,\\
                    \text{(b) hyperbolic:}&\quad
                    \text{infinitely many disjoint lines, angle sum }<\pi,\\
                    \text{(c) elliptic / spherical:}&\quad
                    \text{no disjoint geodesic line, angle sum }>\pi.
                \end{aligned}
            }
        \]
    \end{solution}

    \begin{expectedproblem}
        Match each foundational clue with the corresponding mathematician or
        pair of mathematicians.
        \begin{problemchoices}
            \item diagonal argument and the comparison of infinite
                  cardinalities;
            \item the paradox of the set of all sets that do not contain
                  themselves;
            \item the axiomatic system ZF, with ZFC obtained by adding the
                  Axiom of Choice;
            \item the independence of the Continuum Hypothesis from ZFC,
                  assuming consistency.
        \end{problemchoices}
    \end{expectedproblem}

    \begin{solution}
        The correct associations are
        \[
            \boxed{
                \begin{aligned}
                    \text{(1)}\;&\text{Georg Cantor},\\
                    \text{(2)}\;&\text{Bertrand Russell},\\
                    \text{(3)}\;&\text{Ernst Zermelo and Abraham Fraenkel},\\
                    \text{(4)}\;&\text{Kurt G\"odel and Paul Cohen}.
                \end{aligned}
            }
        \]
    \end{solution}

    \begin{expectedproblem}
        Arrange the following classes of commutative rings from strongest to
        weakest:
        \begin{problemchoices}
            \item unique factorization domain;
            \item field;
            \item integral domain;
            \item principal ideal domain;
            \item Euclidean domain.
        \end{problemchoices}
    \end{expectedproblem}

    \begin{solution}
        The standard implication chain is
        \[
            \boxed{
                \text{field}
                \Longrightarrow
                \text{Euclidean domain}
                \Longrightarrow
                \text{PID}
                \Longrightarrow
                \text{UFD}
                \Longrightarrow
                \text{integral domain}
            }.
        \]
    \end{solution}

    \begin{expectedproblem}
        Identify the Joseon mathematician and chief state councillor who wrote
        \textit{Gusuryak}, exhibited orthogonal Latin squares of order $9$, and
        devised Jisuguimundo, also called the Hexagonal Tortoise Problem.
    \end{expectedproblem}

    \begin{solution}
        The mathematician is
        \[
            \boxed{\text{Choi Seok-jeong}}.
        \]
    \end{solution}

    \begin{expectedproblem}
        Identify the Oxford mathematics lecturer Charles Lutwidge Dodgson by the
        pen name under which he wrote
        \textit{Alice's Adventures in Wonderland}.
    \end{expectedproblem}

    \begin{solution}
        His pen name was
        \[
            \boxed{\text{Lewis Carroll}}.
        \]
    \end{solution}

    \begin{expectedproblem}
        Who received the first Abel Prize, awarded in $2003$?
    \end{expectedproblem}

    \begin{solution}
        The first Abel laureate was
        \[
            \boxed{\text{Jean-Pierre Serre}}.
        \]
    \end{solution}

    \begin{expectedproblem}
        Identify the mathematician associated with the modularity conjecture
        formerly called the Taniyama--Shimura conjecture, a central ingredient
        in the proof of Fermat's Last Theorem.
    \end{expectedproblem}

    \begin{solution}
        The mathematician is
        \[
            \boxed{\text{Goro Shimura}}.
        \]
    \end{solution}

    \begin{expectedproblem}
        What collective pseudonym was adopted by the group of mainly French
        mathematicians who sought a unified axiomatic presentation of modern
        mathematics?
    \end{expectedproblem}

    \begin{solution}
        The collective pseudonym is
        \[
            \boxed{\text{Nicolas Bourbaki}}.
        \]
    \end{solution}

    \begin{expectedproblem}
        Brook Taylor published the general expansion now bearing his name in
        the early eighteenth century, but which Scottish mathematician had
        already found the underlying series method in $1671$?
    \end{expectedproblem}

    \begin{solution}
        The mathematician was
        \[
            \boxed{\text{James Gregory}}.
        \]
    \end{solution}

    \begin{expectedproblem}
        Which French mathematician is remembered for the theorem asserting an
        interior zero of the derivative when a differentiable function has
        equal endpoint values, and also adopted the now-standard indexed
        radical notation for $n$th roots?
    \end{expectedproblem}

    \begin{solution}
        The mathematician is
        \[
            \boxed{\text{Michel Rolle}}.
        \]
    \end{solution}

    \begin{expectedproblem}
        The rule commonly called l'H\^opital's rule appeared in
        l'H\^opital's calculus textbook, but the underlying method came from
        lessons by another mathematician. Identify that mathematician.
    \end{expectedproblem}

    \begin{solution}
        The mathematician was
        \[
            \boxed{\text{Johann Bernoulli}}.
        \]
    \end{solution}

    \begin{expectedproblem}
        In mathematics, the abbreviation FTA commonly refers to two different
        fundamental theorems. Name both.
    \end{expectedproblem}

    \begin{solution}
        They are
        \[
            \boxed{
                \text{the Fundamental Theorem of Algebra and the Fundamental
                Theorem of Arithmetic}
            }.
        \]
    \end{solution}

    \begin{expectedproblem}
        What is the standard name for the ring
        \[
            \Z[i]=\{a+bi\mid a,b\in\Z\}?
        \]
    \end{expectedproblem}

    \begin{solution}
        Its elements are the
        \[
            \boxed{\text{Gaussian integers}}.
        \]
    \end{solution}

    \begin{expectedproblem}
        The standard mathematical model of single-fold origami is usually
        described by seven axioms named after two mathematicians. What are
        these axioms commonly called?
    \end{expectedproblem}

    \begin{solution}
        They are the
        \[
            \boxed{\text{Huzita--Hatori axioms}}.
        \]
        The historical naming is not a strict priority claim: Jacques Justin
        had published all seven axioms earlier, while Humiaki Huzita and
        Koshiro Hatori are associated with the commonly used name.
    \end{solution}

    \begin{expectedproblem}
        What branch of mathematics studies objects together with morphisms
        between them, with associative composition and identity morphisms?
    \end{expectedproblem}

    \begin{solution}
        The branch is
        \[
            \boxed{\text{category theory}}.
        \]
    \end{solution}

    \begin{expectedproblem}
        In the title of al-Khwarizmi's foundational algebra treatise, two Arabic
        terms describe restoring a subtracted quantity and balancing or
        reducing corresponding quantities. Name the two terms.
    \end{expectedproblem}

    \begin{solution}
        The terms are
        \[
            \boxed{\textit{al-jabr}\quad\text{and}\quad\textit{al-muqabala}}.
        \]
        The word \emph{algebra} derives from \emph{al-jabr}.
    \end{solution}

    \begin{expectedproblem}
        How many reachable configurations does the standard $3\times3\times3$
        Rubik's Cube have?
    \end{expectedproblem}

    \begin{solution}
        The corner pieces can be permuted in $8!$ ways and oriented in $3^7$
        independent ways. The edges can be permuted in $12!$ ways and oriented
        in $2^{11}$ ways, but the corner and edge permutation parities must
        agree, giving a factor of $1/2$. Thus
        \[
            8!\,3^7\,12!\,2^{10}
            =
            \boxed{43\,252\,003\,274\,489\,856\,000}.
        \]
    \end{solution}

    \begin{expectedproblem}
        For a regular space curve, what is the standard collective name for the
        data
        \[
            \{\mathbf{T},\mathbf{N},\mathbf{B},\kappa,\tau\},
        \]
        consisting of the unit tangent, principal normal, binormal, curvature,
        and torsion?
    \end{expectedproblem}

    \begin{solution}
        This collection is the
        \[
            \boxed{\text{Frenet--Serret apparatus}}.
        \]
    \end{solution}

    \begin{expectedproblem}
        The median identity
        \[
            AB^2+AC^2
            =
            2\bigl(AM^2+BM^2\bigr),
        \]
        where $M$ is the midpoint of $BC$, is sometimes called the Pappus
        median theorem in the Republic of Korea and Japan. What is its standard international
        name?
    \end{expectedproblem}

    \begin{solution}
        It is
        \[
            \boxed{\text{Apollonius' theorem}}.
        \]
    \end{solution}

    \begin{expectedproblem}
        A fifteenth family of convex pentagons capable of tiling the plane was
        found in $2015$. By $2017$, an exhaustive computer-assisted search had
        settled the classification. As of $2026$, how many families of convex
        pentagons that tile the plane are known to exist in this classification?
    \end{expectedproblem}

    \begin{solution}
        The exhaustive classification found no additional families beyond the
        previously known fifteen. Therefore the number is
        \[
            \boxed{15}.
        \]
    \end{solution}

